% Text recovered from pages 57--59 of the original complete reader PDF; mathematical notation flattened by extraction.
The problem of the chapter. Grothendieck starts from a category of coefficients, here the category of small categories itself, and asks how its homotopy theory should vary with a diagram shape. The value at a small category S is to be


HOTW (S) = Fun(Sop, Cat)[W \ensuremath{-}1


obj], where a natural transformation belongs to Wobj when all its components belong to W. The essential operation is to localize the category of diagrams before forgetting their coherence. A functor from Sop to the already localized category of coefficients generally retains less information. This is the same motivation that led to the diagram formulation in Grothendieck (1983, \S{}69), and that the letter to Thomason subsequently places at the centre of the derivator programme (Grothendieck, 1991). The discussions of coherence and evaluation in the commentary on Chapter~I, manuscript pages 7, 17--18 and 27, remain the background; the present chapter supplies categorical constructions rather than repeating those objections.


The first seventy manuscript pages, dated 14--15 November 1990, record the discovery of the correct comparison theorem. Thirteen categories of objects over S initially appear to offer interchangeable models. They include split fibrations, arbitrary fibrations, smooth and proper functors, homotopy fibrations, perfect functors, and all of Cat/S. The objects alone do not decide the homotopy theory. The permitted morphisms and the chosen weak equivalences matter just as much. The first theorem overlooks that distinction, is expressly withdrawn, and is replaced by a comparison between three theories. The manuscript’s revisions are part of the proof’s development, not noise to remove from it. The axioms and the four tests. The axiom list also changes. The original L4 is an absolute Theorem A: an aspherical functor is a weak equivalence. L4 bis, introduced on manuscript page 44, is the relative assertion about a map X \ensuremath{\to}Y over a third category S, tested on every map X/s \ensuremath{\to}Y/s. It is the latter formulation that supports the relative descent arguments. The initial L5 passes from fibrewise weak equivalences of Cartesian maps between fibrations to weak equivalences of total categories. The new L5 on page 43 goes in the opposite direction for maps between homotopy fibrations. Their common label must not obscure their different content. The commentary proves the initial claim that L1, L2 bis and the original L5 already imply the required elementary axioms, including relative descent, by a twisted-arrow argument for opposite invariance.


For a map f : X \ensuremath{\to}Y over S, the decisive tests are


coarse


f \ensuremath{\in}W,


fibrewise fs : Xs \ensuremath{\to}Ys belongs to W for every s,


left-fine


f/s : X/s \ensuremath{\to}Y/s belongs to W for every s,


right-fine


s\textbackslash{}f : s\textbackslash{}X \ensuremath{\to}s\textbackslash{}Y belongs to W for every s. Over a smooth object, the strict fibre compares to the undercategory; over a proper object, it compares to the overcategory. These comparisons and relative descent explain when fibrewise equivalence implies a universal equivalence of totals. The modern asphericity formulation retains precisely this mechanism (Maltsiniotis, 2005, Theorem 1.2, Corollary 4.16 and Theorem 5.4). Coarse equivalence alone does not give either fine condition on the external borders of the diagram. Over [1], the inclusion \{0\} \ensuremath{\to}[1] is a coarse equivalence of split fibrations, but its fibre at 1 is the map from the empty category to a point. Evaluation at that fibre proves that the two localizations themselves differ. The two comma resolutions. The functor \ensuremath{\Phi} replaces X \ensuremath{\to}S by the split fibration whose fibre at s is s\textbackslash{}X; \ensuremath{\Psi} gives the split opfibration with fibre X/s. Their units are categorical homotopy equivalences on total categories. More precisely, the unit of \ensuremath{\Phi} is right-fine, and the unit of \ensuremath{\Psi} is left-fine. Composition of the two arrows in a comma object gives an explicit contraction proving this stronger fact. Consequently the right-fine localization of Cat/S computes the contravariant diagram theory, while the left-fine localization computes its covariant counterpart.


The comparison with chosen splittings needs a genuine repair of the argument, but not a change in the desired theorem. A Cartesian functor need not preserve the selected Cartesian lifts strictly, and an arbitrary functor over the base need not preserve Cartesian arrows at all. The actual free split adjunction is


Homsplit/S(\ensuremath{\Phi}X, Y ) \ensuremath{\sim}= HomS(X, UY ).


Its counit supplies the comparison in the strict category. For arbitrary functors between fibrations the pullback retraction has a canonical, possibly noninvertible natural comparison. Natural transformations over S induce equality after fine localization, since the cylinder projection A \ensuremath{\times} [1] \ensuremath{\to}A is fine. This yields the natural isomorphism needed after localization. The exposition constructs the adjunction and the localized comparison in full; the historical and modern strictification references are separated from this particular elaboration (Gray, 1966; Emmenegger et al., 2024).


The covariant theory is HOTW (Sop). The additional outer opposite that appears in several passages of the edition is incorrect with the initial definition above: taking opposite total categories preserves the direction of functors between them. The corrected variance is used throughout the commentary, while both the original statements and the editorial brackets are preserved in the source. The strength of path completion. The two-sided path replacement gives a natural map X \ensuremath{\to}C(X) which is a weak homotopy equivalence, with C(X) \ensuremath{\to}S perfect. This immediately proves comparisons for the coarse localizations, without the new L5. To identify coarse and fibrewise equivalences on the internal categories, and then identify their common theory with a full subcategory of the fine theory, is stronger. That step is what the new L5 supplies.


Its restriction on W is severe. A degree-n covering of the circle is an equivalence for homology with coefficients in a ring where n is invertible, but its fibre has n components. The commentary realizes the covering as an action groupoid, so that perfectness and the strict fibre can be checked


directly. For integral homology it gives a different explicit example:


Z[1/2] \ensuremath{\rtimes}2 Z


  \ensuremath{-}\ensuremath{\to}BZ. B The fibre has first homology Z[1/2]. The mapping-torus sequence shows that the total projection is nevertheless an integral homology equivalence, because 1 \ensuremath{-}2 acts invertibly on that fibre homology. This calculation settles the tentative integral-coefficient assertion on manuscript page 43. The development label marks the supplied application and proof, with a bounded account of the literature checked; it claims no priority for the group or the homology calculation.


Finite homotopy truncations fail as well: a path fibration over K(A, n + 1) has contractible total space and fibre K(A, n). Retaining the standing condition W\ensuremath{\infty}\ensuremath{\subseteq}W \ensuremath{\subseteq}W0, the new L5 actually forces W = W\ensuremath{\infty}. Iterating path fibres shows that every W-aspherical category has connected iterated loop spaces, hence is weakly contractible; applying this to the fibres of a replacement of an arbitrary W-map proves the conclusion. The shortened annotation omitting the component condition cannot stand literally, since the maximal class of all functors is a counterexample. Cisinski’s later characterization explicitly includes the corresponding condition on components (Cisinski, 2004, Theorem 2.3.6). The reflection and its modern meaning. The final argument proves full faithfulness through an adjunction. Put


\ensuremath{\Sigma}0 = W d


A = Cat/S, S \ensuremath{\subseteq}\ensuremath{\Sigma} = WS. Path completion carries \ensuremath{\Sigma} into \ensuremath{\Sigma}0, its unit \ensuremath{\eta}X belongs to \ensuremath{\Sigma}, and \ensuremath{\eta}CX belongs to \ensuremath{\Sigma}0. These conditions give


P : A[\ensuremath{\Sigma}\ensuremath{-}1


0 ] \ensuremath{\rightleftarrows}A[\ensuremath{\Sigma}\ensuremath{-}1] : R, P \ensuremath{\dashv}R, with R fully faithful. The unit is induced by \ensuremath{\eta}, and the counit is its inverse in the coarse category. The commentary verifies both triangle identities. In particular, the equality C(\ensuremath{\eta}X) = \ensuremath{\eta}CX after fine localization follows formally from naturality and the stated invertibilities. Grothendieck’s concatenation of paths gives an additional geometric proof by providing a common inverse. There is no need to cancel \ensuremath{\eta}X in the fine category, where it need not be invertible.


For W = W\ensuremath{\infty}, the resulting interpretation is


HOTW\ensuremath{\infty}(S) \ensuremath{\simeq}Ho Fun(N(S)op, S),


HOTlc


W\ensuremath{\infty}(S) \ensuremath{\simeq}Ho Fun(\ensuremath{\Pi}\ensuremath{\infty}N(S), S) \ensuremath{\simeq}Ho(S/|N(S)|). Here S is the infinity-category of spaces, and the homotopy category is taken after forming coherent diagrams or the slice. The reflector is restriction after left Kan extension along groupoid completion. The commentary explains the comparison through Thomason’s model and straightening (Thomason, 1980; Lurie, 2009). It also computes the interval reflector and the endomorphisms of the point over BZ, which distinguish the homotopy category of the slice from the ordinary slice of the homotopy category.


The locally constant embedding and its two adjoints are now theorems in the derivator setting (Cisinski, 2009, Proposition 3.10 and Theorem 3.13). Their existence does not identify an arbitrary coarse W-localization with that subcategory. The surviving qualification is about hypotheses, not an unresolved version of the full-faithfulness problem for spaces. We also do not label the


relationship between the weakest absolute and relative axiom packages as a current open problem: this audit establishes the implications used here and records where a weaker package is insufficient for the printed proof. The construction of the path categories themselves is an input from outside these pages. The first part of the chapter ends in the opening clauses of the theorem--scholium on manuscript page 70; the continuation on pages 71--81 is commented in the second part.

\nocite{*}
